\documentclass[a4paper]{article}

\usepackage{graphicx}
\usepackage{amsmath,amssymb}
\usepackage{minted}
\usepackage{multirow}
\usepackage{float}
\usepackage{algorithm}
\usepackage{algpseudocode}
\usepackage{todonotes}
\usepackage{forest}
\usepackage{xstring}
\usepackage{xcolor}
\usepackage[most]{tcolorbox}
\usepackage{xparse}
\usepackage{natbib}
\usepackage{adjustbox}

\usetikzlibrary{graphs,graphdrawing}
\usegdlibrary{layered}

% for external graphviz
\input{/home/donkarlo/Dropbox/repo/nd_language_ai_project/src/nd_language_ai/formal/computer/programming/tex/latex/code_clips/external_graphviz.tex}

% for tags
\input{/home/donkarlo/Dropbox/repo/nd_language_ai_project/src/nd_language_ai/formal/computer/programming/tex/latex/part/control_sequence/kind/semtag/semtag.tex}

% for links
\input{/home/donkarlo/Dropbox/repo/nd_language_ai_project/src/nd_language_ai/formal/computer/programming/tex/latex/code_clips/seehere.tex}

% for nested lists and sections
\input{/home/donkarlo/Dropbox/repo/nd_language_ai_project/src/nd_language_ai/formal/computer/programming/tex/latex/code_clips/deep_structure.tex}

\title{Artificial Potential Field}
\date{}
\author{Mohammad Rahmani}

\begin{document}
   \maketitle

   \section{Definition}
      The artificial potential field (APF) method represents a navigation problem by assigning a virtual potential to the robot configuration. The goal produces an attractive contribution and obstacles produce repulsive contributions. The robot is then driven in the direction of decreasing total potential.

      Khatib introduced the artificial potential field formulation for real-time obstacle avoidance in manipulators and mobile robots \cite{khatib-1986-real-time-obstacle-avoidance}.

   \section{Potential and virtual force}
      Let \(\mathbf q\) denote the robot position or configuration. The total potential is
      \[
         U(\mathbf q)
         =
         U_{\mathrm{att}}(\mathbf q)
         +
         U_{\mathrm{rep}}(\mathbf q).
      \]

      The corresponding virtual force or guidance vector is the negative gradient of the potential,
      \[
         \mathbf F(\mathbf q)
         =
         -\nabla U(\mathbf q)
         =
         \mathbf F_{\mathrm{att}}(\mathbf q)
         +
         \mathbf F_{\mathrm{rep}}(\mathbf q).
      \]

   \section{Attractive field}
      A common attractive potential around a goal \(\mathbf q_g\) is
      \[
         U_{\mathrm{att}}(\mathbf q)
         =
         \frac{1}{2}\zeta
         \left\|
            \mathbf q-\mathbf q_g
         \right\|^2,
      \]
      where \(\zeta>0\) controls the strength of attraction. Its negative gradient is
      \[
         \mathbf F_{\mathrm{att}}(\mathbf q)
         =
         -\zeta
         \left(
            \mathbf q-\mathbf q_g
         \right).
      \]

   \section{Repulsive field}
      Let \(\rho(\mathbf q)\) be the distance from the robot to the nearest obstacle, and let \(\rho_0\) be the obstacle influence distance. A common repulsive potential is
      \[
         U_{\mathrm{rep}}(\mathbf q)
         =
         \begin{cases}
            \dfrac{1}{2}\eta
            \left(
               \dfrac{1}{\rho(\mathbf q)}
               -
               \dfrac{1}{\rho_0}
            \right)^2,
            & \rho(\mathbf q)\leq \rho_0,\\[8pt]
            0,
            & \rho(\mathbf q)>\rho_0,
         \end{cases}
      \]
      where \(\eta>0\) controls the repulsive strength. The resulting vector points away from the obstacle because it follows the negative gradient of the repulsive potential.

   \section{Relation to the control input}
      The APF output should not automatically be interpreted as a physical force. It is a guidance vector. Depending on the robot model and controller, it can be mapped to a desired force, acceleration, velocity, heading, or another control reference. For example,
      \[
         \mathbf u_t
         =
         k\,\mathbf F(\mathbf q_t)
      \]
      can be used when the control-input semantics are chosen so that \(\mathbf u_t\) represents the corresponding command.

   \section{Biological analogue}
      Artificial potential fields do not have a literal one-to-one biological equivalent. A close biological analogue is gradient-following behavior such as chemotaxis. Bacteria can bias their movement toward attractants and away from repellents by sensing environmental gradients \cite{hu-2014-bacterial-navigation-gradients}. This resembles the attractive and repulsive components of an APF at the behavioral level, although the biological mechanism is implemented through sensing, signaling, adaptation, and stochastic run-and-tumble motion rather than by explicitly computing a scalar artificial potential.

      \cite{kerr-2017-biological-goal-seeking} also uses biologically inspired goal-seeking mechanisms in autonomous agents and combines potential-field-like behavior with finite-state control.

   \section{Limitation}
      A classical limitation of artificial potential fields is the possibility of local minima: the total gradient can become zero at a configuration that is not the goal. Therefore, APF is particularly suitable as a reactive or local navigation mechanism, while additional strategies may be required for global completeness in complex environments.

   \bibliographystyle{plainnat}
   \bibliography{/home/donkarlo/Dropbox/repo/nd_shared_research/bibliography/bibliography.bib}
\end{document}
